\noindent{\large\bfseries Practice questions}\hfill 28 October 2026\par
\vspace{2pt}
\noindent\rule{\linewidth}{0.6pt}\par
\vspace{2pt}
{\small Use $\pi=\dfrac{22}{7}$. For every question without a picture, first draw a rough diagram and mark all the given measurements on it. Show all your steps.\par}
\ifteacher
\vspace{4pt}
{\color{thinrule}\noindent\rule{\linewidth}{0.4pt}}\par
\vspace{4pt}
{\small\color{muted}\linespread{1.15}\selectfont
\textbf{Purpose.} This set revises the whole chapter. The questions follow the style of the NCERT textbook \textit{Ganita Manjari}, Grade~9, Chapter~6. They go from easier to harder.\nl
\textbf{Competency focus.} Students work backwards from a result, find and explain a mistake, test a claim, compare two methods, and apply results to real situations. Several questions ask students to draw a diagram.\nl
\textbf{Concepts assessed.} Circumference of a circle (Q1, 12); the number $\pi$ (Q2); arc length (Q3, 13); the stagger on a running track (Q4); area of a parallelogram (Q5, 7); area of a triangle and triangles between the same parallels (Q6); the median of a triangle (Q7); Heron's formula (Q8, 9); Brahmagupta's formula (Q10); Baudh\={a}yana's method of squaring a rectangle (Q11); area of a circle (Q12, 15); area of a sector and a ring sector (Q13); area of a segment (Q14); composite figures (Q15).\nl
\textbf{Answers.} Answers appear in grey after each question, with a rough answer diagram.\par}
\vspace{4pt}
{\color{thinrule}\noindent\rule{\linewidth}{0.4pt}}\par
\vspace{3pt}
\fi
